% ----------------------------------------------------------
% Conclusão
% ----------------------------------------------------------
\chapter[Conclusão]{Conclusão}
\label{cap:conclusao}
%\addcontentsline{toc}{chapter}{Conclusão}
% ---

Este trabalho teve como objetivo desenvolver uma ferramenta para auxiliar os discentes dos cursos da área de computação no aprendizado de linguagens formais e máquinas de estados. O foco principal recaiu sobre os autômatos de pilha não determinísticos (APNs), visto que esses conhecimentos, embora fundamentais, caracterizam-se por sua alta complexidade e abstração.

Entre as principais contribuições do projeto desenvolvido, destacam-se: a modelagem gráfica de APNs; a visualização, passo a passo, das múltiplas ramificações de execução do autômato; e a funcionalidade de criação e correção de exercícios. Todos esses recursos foram integrados em uma plataforma web, eliminando qualquer necessidade de instalação prévia no computador do usuário.


Apesar dos resultados satisfatórios, o sistema apresenta algumas limitações. A principal delas diz respeito à execução de autômatos que podem gerar um número exponencial de caminhos ao computarem palavras muito extensas. Como a aplicação é executada integralmente no lado do cliente (\textit{frontend}), essa restrição varia significativamente de acordo com a capacidade do \textit{hardware} do usuário. No entanto, essa limitação não representa um obstáculo significativo, uma vez que a ferramenta foi concebida com propósitos didáticos para a simulação de exemplos acadêmicos, e não para o processamento de modelos de grande porte voltados à pesquisa.

Dessa forma, conclui-se que a ferramenta desenvolvida demonstra grande potencial para auxiliar os alunos na compreensão dos autômatos de pilha não determinísticos. Contudo, ressalta-se que ainda não foi realizada uma pesquisa prática com o público-alvo para determinar empiricamente sua eficácia e efetividade. Sendo assim, a validação com os usuários consolida-se como um dos principais tópicos para trabalhos futuros. Adicionalmente, sugere-se a implementação de outras máquinas de estados além dos APNs, bem como o desenvolvimento de um \textit{backend} para o sistema, o que otimizaria o desempenho da aplicação e permitiria a execução de autômatos mais complexos.
